\documentclass[11pt]{article}
\usepackage[a4paper,margin=1in]{geometry}
\usepackage{amsmath,amssymb,amsthm,mathtools}
\usepackage{booktabs}
\usepackage{microtype}
\usepackage[hidelinks]{hyperref}
\hypersetup{
  pdftitle={An Explicit Infinite Model Refuting Ulrich's u4 as a Single Axiom for Positive Implication},
  pdfauthor={Byungwoong Yoo},
  pdfsubject={Machine-verified countermodel for Ulrich's u4},
  pdfkeywords={positive implication, single axiom, Lean 4, countermodel, Ulrich u4}
}

\newtheorem{theorem}{Theorem}
\newtheorem{lemma}[theorem]{Lemma}
\newtheorem{corollary}[theorem]{Corollary}
\newcommand{\Imp}{\mathbin{\to}}
\newcommand{\Inst}{\operatorname{Inst}}

\title{An Explicit Infinite Model Refuting Ulrich's \texttt{u4}\\
as a Single Axiom for Positive Implication}
\author{Byungwoong Yoo\\Independent Researcher}
\date{Corrected strengthened release v2\\
Lean 4.33.0, leanchecker, nanoda, axiom-audit, and symbolic-check suite\\
Proof-source commit \texttt{fec164ed763794330d57f76ddd1b0c390f22db09}\\
\href{https://github.com/ByungwoongYoo/ntg-optn-tbk1-tkm-triage/actions/runs/32384585278}
{All-green GitHub Actions run 32384585278}\\
Verification snapshot: 20 August 2026 UTC (21 August 2026 KST)\\
Reserved corrected-v2 DOI: \href{https://doi.org/10.5281/zenodo.22031656}
{\texttt{10.5281/zenodo.22031656}}\\
Corrected v2 publication: 21 August 2026 KST}

\begin{document}
\maketitle

\begin{abstract}
Fitelson and Peltier (2026) eliminated three of Ulrich's four remaining
15-symbol candidates for a single axiom of positive implicational logic,
leaving
\[
((x\Imp y)\Imp z)\Imp((y\Imp(z\Imp u))\Imp(y\Imp u))
\]
as an open case.  We construct an explicit infinite model satisfying every
substitution instance of this formula and closed under modus ponens, while
falsifying reflexivity $q\Imp q$.  The model is the union of the substitution
instances of a recursively defined family of formula schemes.  The decisive
countermodel property, uniform-substitution closure, the induced proof system,
and an internal separation from the standard positive-implication basis are
formalized in Lean 4.33.0.  At the proof-source commit identified above, the
Lean kernel, \texttt{leanchecker}, the independently implemented Rust checker
\texttt{nanoda}, a compiled-environment axiom audit, and a separate Python
symbolic checker all succeeded.  These are machine-verification results within
the project; external specialist and journal review remain pending.
\end{abstract}

\section{Background}

Carew Meredith found single axioms of length 17 for positive implication.
Dolph Ulrich reduced the possibility of a shorter axiom to four formulas of
length 15.  Fitelson and Peltier used saturation-based theorem proving to
eliminate the first three and reported that the fourth candidate,
\begin{equation}\label{eq:u4}
\Phi=((x\Imp y)\Imp z)\Imp((y\Imp(z\Imp u))\Imp(y\Imp u)),
\end{equation}
remained open.  They further established that the shortest single axioms
have length 17 if and only if $\Phi$ is not a single axiom.

The proof below is semantic.  Its domain is the free algebra of finite
implicational formulas.  The theorem predicate is a particular substitution-
closed set that is also closed under modus ponens.

\section{Recursive schemes}

Write
\[
C(A,B,R):=(A\Imp(B\Imp R))\Imp(A\Imp R),
\]
so that
\[
\Phi=((x\Imp y)\Imp z)\Imp C(y,z,u).
\]

Choose pairwise distinct variables
\[
a,b,c,r_2,r_3,\ldots
\]
and define
\[
T_1=\Phi,
\]
\[
A_2=a,\qquad B_2=C(b,a,c),\qquad T_2=C(A_2,B_2,r_2).
\]
For $n\ge2$, recursively set
\[
A_{n+1}=B_n\Imp r_n,\qquad
B_{n+1}=A_n\Imp r_n,
\]
\[
T_{n+1}=C(A_{n+1},B_{n+1},r_{n+1}).
\]
For $n\ge2$, put
\[
D_n=A_n\Imp(B_n\Imp r_n),
\]
so that
\[
T_n=D_n\Imp(A_n\Imp r_n).
\]

For a formula scheme $S$, let $\Inst(S)$ denote the set of all finite
substitution instances of $S$.

\section{\hspace{0.35em}The interlocking property}

For finite formulas $P,Q$, let $E(P,Q)$ mean that no substitution $\sigma$
and finite formula $W$ satisfy
\[
\sigma(P)=\sigma(Q)\Imp W.
\]

\begin{lemma}[Interlocking]\label{lem:interlocking}
For every $n\ge2$,
\[
E(A_n,B_n)\qquad\text{and}\qquad E(B_n,A_n).
\]
\end{lemma}

\begin{proof}
At $n=2$,
\[
A_2=a,\qquad
B_2=((b\Imp(a\Imp c))\Imp(b\Imp c)).
\]
If $\sigma(a)=\sigma(B_2)\Imp W$, then the finite formula
$\sigma(a)$ occurs as a proper subformula of its own right-hand side.
This is impossible.  Conversely, if $\sigma(B_2)=\sigma(a)\Imp W$, injectivity
of implication forces
\[
\sigma(b\Imp(a\Imp c))=\sigma(a),
\]
which again makes $\sigma(a)$ a proper subformula of itself.

Assume both assertions at $n$.  Since
\[
A_{n+1}=B_n\Imp r_n,\qquad
B_{n+1}=A_n\Imp r_n,
\]
an equality
\[
\sigma(A_{n+1})=\sigma(B_{n+1})\Imp W
\]
would imply
\[
\sigma(B_n)=\sigma(A_n)\Imp\sigma(r_n),
\]
contradicting $E(B_n,A_n)$.  The converse direction is symmetric.
\end{proof}

\section{The major-premise obstruction}

\begin{lemma}\label{lem:major}
For every $m\ge2$ and $n\ge1$, the schemes $D_m$ and $T_n$
have no common substitution instance.
\end{lemma}

\begin{proof}
First let $n\ge2$, and suppose
\[
\sigma(D_m)=\tau(T_n).
\]
Expanding both sides and using injectivity of implication gives
\[
\sigma(A_m)=\tau(A_n)\Imp(\tau(B_n)\Imp\tau(r_n))
\]
and
\[
\sigma(B_m)=\tau(A_n).
\]
Thus
\[
\sigma(A_m)=\sigma(B_m)\Imp(\tau(B_n)\Imp\tau(r_n)),
\]
contrary to Lemma~\ref{lem:interlocking}.

Now let $n=1$.  If $m\ge3$, then
\[
A_m=B_{m-1}\Imp r_{m-1},\qquad
B_m=A_{m-1}\Imp r_{m-1}.
\]
Matching $D_m$ with $T_1$ forces both
\[
\sigma(r_{m-1})=\tau(z)
\quad\text{and}\quad
\sigma(r_{m-1})=\tau(z\Imp u),
\]
hence $\tau(z)=\tau(z)\Imp\tau(u)$, impossible for a finite formula.

For $m=2$, a common instance of
\[
D_2=a\Imp(C(b,a,c)\Imp r_2)
\]
and $T_1$ would imply
\[
\sigma(a)=\tau((x\Imp y)\Imp z)
\]
and
\[
\sigma(C(b,a,c))=\tau(y\Imp(z\Imp u)).
\]
The second equality yields
\[
\tau(y)=\sigma(b\Imp(a\Imp c)),\quad
\tau(z)=\sigma(b),\quad
\tau(u)=\sigma(c).
\]
Substitution into the first equality makes $\sigma(a)$ contain itself
as a proper subformula, a contradiction.
\end{proof}

\section{The infinite model}

Define
\[
\mathcal P=\bigcup_{n\ge1}\Inst(T_n).
\]
Interpret the unary theorem predicate $p$ on the free algebra of finite
implicational formulas by
\[
p(F)\quad\Longleftrightarrow\quad F\in\mathcal P.
\]

Every instance of $\Phi=T_1$ belongs to $\mathcal P$.  It remains to prove
closure under modus ponens.

\begin{lemma}\label{lem:mp}
If $X\in\mathcal P$ and $X\Imp Y\in\mathcal P$, then $Y\in\mathcal P$.
\end{lemma}

\begin{proof}
Choose $n,m\ge1$ and substitutions $\alpha,\beta$ with
\[
X=\alpha(T_n),\qquad X\Imp Y=\beta(T_m).
\]
If $m\ge2$, the antecedent of $\beta(T_m)$ is $\beta(D_m)$.
Then $X$ would be a common instance of $D_m$ and $T_n$, contrary to
Lemma~\ref{lem:major}.  Hence $m=1$, so
\[
X=\beta((x\Imp y)\Imp z),\qquad Y=\beta(C(y,z,u)).
\]

If $n=1$, rename the variables of the minor copy of $T_1$ as
$x',y',z',u'$.  Comparison of
\[
\alpha(T_1)=\beta((x\Imp y)\Imp z)
\]
gives
\[
\beta(y)=\alpha(z'),\qquad
\beta(z)=\alpha(C(y',z',u')).
\]
Therefore
\[
Y=C\bigl(\alpha(z'),\alpha(C(y',z',u')),\beta(u)\bigr),
\]
an instance of $T_2=C(a,C(b,a,c),r_2)$.

If $n\ge2$, comparison of
\[
\alpha(T_n)=\beta((x\Imp y)\Imp z)
\]
gives
\[
\beta(y)=\alpha(B_n\Imp r_n),\qquad
\beta(z)=\alpha(A_n\Imp r_n).
\]
Hence
\[
Y=C\bigl(\alpha(B_n\Imp r_n),
         \alpha(A_n\Imp r_n),\beta(u)\bigr),
\]
which is an instance of $T_{n+1}$.  Thus $Y\in\mathcal P$.
\end{proof}

\section{\hspace{0.35em}Failure of reflexivity}

\begin{lemma}\label{lem:refl}
No formula $Q\Imp Q$ belongs to $\mathcal P$.
\end{lemma}

\begin{proof}
For $n\ge2$, an instance of
\[
T_n=(A_n\Imp(B_n\Imp r_n))\Imp(A_n\Imp r_n)
\]
could equal $Q\Imp Q$ only if
\[
\sigma(B_n\Imp r_n)=\sigma(r_n).
\]
This would make the finite formula $\sigma(r_n)$ contain itself as a
proper subformula.

For $n=1$, equality of the antecedent and consequent of an instance of
$T_1$ forces
\[
\sigma(y)=\sigma(z\Imp u),\qquad
\sigma(z)=\sigma(y\Imp u),
\]
a contradiction by strict size increase in both directions.
\end{proof}

\begin{theorem}\label{thm:main}
Ulrich's formula \eqref{eq:u4} is not a single axiom for positive
implicational logic under substitution and modus ponens.
\end{theorem}

\begin{proof}
The interpretation above satisfies every substitution instance of $\Phi$,
and Lemma~\ref{lem:mp} proves that it satisfies modus ponens.  Lemma~
\ref{lem:refl} shows that it falsifies reflexivity $Q\Imp Q$, which is a
theorem of positive implication.  Therefore reflexivity is not derivable
from $\Phi$, so $\Phi$ is not a single axiom.
\end{proof}

\begin{corollary}
Combining Theorem~\ref{thm:main} with the elimination and classification
results established by Fitelson and Peltier, the minimum length of a single
axiom for positive implication is 17 symbols.
\end{corollary}

\section{Formal certificate and verified release}

The evidence bundle includes \texttt{lean/U4Formal.lean}.  It defines
the free algebra of finite implicational formulas, the exact formula $\Phi$,
and an explicit predicate $P$.  The Lean kernel-checked theorem is
\begin{verbatim}
theorem U4Formal.u4_countermodel_exists :
    Exists fun pred : Fml -> Prop =>
      (forall x y z u, pred (Phi x y z u)) /\
      (forall X Y, pred X -> pred (I X Y) -> pred Y) /\
      (forall q, Not (pred (I q q)))
\end{verbatim}
The strengthened file additionally formalizes arbitrary uniform substitution,
proves that every syntactic \texttt{u4} derivation is sound in the model, and
derives reflexivity from the standard $K$ and $S$ schemes.  Its final theorem
exhibits a positive-implication theorem that is not derivable from \texttt{u4}:
\begin{verbatim}
theorem U4Formal.u4_not_axiomatizes_positive_implication :
    Exists fun f : Fml => PositiveDerivable f /\ Not (U4Derivable f)
\end{verbatim}
The source scan forbids \texttt{sorry}, \texttt{admit}, custom
\texttt{axiom}, and \texttt{unsafe}; the compiled-environment audit is the
check of the declared axiom-dependency boundary.  The CI artifact records the exact
Git commit, all checker outcomes, direct compiler output, source files, and
SHA-256 hashes.  Run 32384585278 passed every gate at commit
\texttt{fec164ed763794330d57f76ddd1b0c390f22db09}.  The same run reports that
\texttt{nanoda} checked 96,602 declarations with no typechecker errors; its one
reported diagnostic concerned pretty-printing axioms and not type checking.

\section{Executable cross-checks}

The release workflow performs the following mutually reinforcing checks.

\begin{center}
\begin{tabular}{@{}lll@{}}
\toprule
Check & Range & Result\\
\midrule
Lean build and direct compilation & Lean 4.33.0 & passed\\
Lean environment recheck & \texttt{leanchecker} & passed\\
Separately implemented Rust type check & \texttt{nanoda}, no \texttt{sorryAx} use & passed\\
Compiled axiom audit & standard allowlist only & passed\\
Python recurrence/interlocking/major obstruction & $T_1,\ldots,T_{100}$ & passed\\
Python exhaustive binary-tree enumeration & through 12 leaves & passed\\
\bottomrule
\end{tabular}
\end{center}

At 12 leaves there are $C_{11}=58{,}786$ binary application-tree shapes.
Exactly one was typable: the right comb.  These checks support, but do not
replace, the all-$n$ proof above.

\section{Verification status and claim boundary}

As of 21 August 2026, the official 2026 source cited below reported
\texttt{u4} as open.  This release makes no claim that all later, unpublished,
or non-indexed work has been identified.

The public identifiers are:
\begin{center}
\begin{tabular}{@{}ll@{}}
Corrected v2 (reserved) &
\href{https://doi.org/10.5281/zenodo.22031656}{\texttt{10.5281/zenodo.22031656}}\\
All-version concept &
\href{https://doi.org/10.5281/zenodo.21987698}{\texttt{10.5281/zenodo.21987698}}\\
Historical v1 &
\href{https://doi.org/10.5281/zenodo.21987699}{\texttt{10.5281/zenodo.21987699}}\\
\end{tabular}
\end{center}
The v1 record contains a Lean 4.30.0 kernel check of the core countermodel.
This corrected
strengthened release adds explicit substitution and derivability, an internal
$K/S$ separation theorem, and multiple independently implemented compiled
checks.  The mathematical exposition, the connection to the prior reduction,
external specialist review, journal peer review, and worldwide priority remain
separate questions.  The defensible description is:

\begin{quote}
Lean-checked explicit infinite countermodel, supported by project-run checks
using separately implemented environment and type checkers, showing that
Ulrich's \texttt{u4} is not a single axiom.  No independent third-party
specialist or journal review has yet occurred, and comparative priority review
remains pending.
\end{quote}

\paragraph{AI-use disclosure.}
OpenAI GPT-5-class systems, including Codex, substantially assisted with proof
exploration, adversarial checking, program drafting, literature-search
assistance, and manuscript preparation.  The human author, Byungwoong Yoo, finalized the
mathematical argument and the submitted package, reviewed the claim and
evidence, and assumes full responsibility for the accuracy and final wording.
Verification and replay details are documented separately in the supplied logs
and checkers.

\begin{thebibliography}{9}

\bibitem{FitelsonPeltier2026}
B. Fitelson and N. Peltier,
``Applying Saturation-Based Theorem Proving to Open Problems in Positive
Implicational Logic,''
\emph{Journal of Automated Reasoning}, vol. 70, article 5, 2026.
\href{https://doi.org/10.1007/s10817-026-09752-1}
{doi:10.1007/s10817-026-09752-1}.

\bibitem{HindleyMeredith1990}
J. R. Hindley and D. Meredith,
``Principal Type-Schemes and Condensed Detachment,''
\emph{Journal of Symbolic Logic}, vol. 55, no. 1, pp. 90--105, 1990.
\href{https://doi.org/10.2307/2274956}{doi:10.2307/2274956}.

\bibitem{Kalman1983}
J. A. Kalman,
``Condensed Detachment as a Rule of Inference,''
\emph{Studia Logica}, vol. 42, no. 4, pp. 443--451, 1983.
\href{https://doi.org/10.1007/BF01371632}{doi:10.1007/BF01371632}.

\bibitem{Ulrich1999}
D. Ulrich,
``New Single Axioms for Positive Implication,''
\emph{Bulletin of the Section of Logic}, vol. 28, pp. 39--42, 1999.

\bibitem{Ulrich2001}
D. Ulrich,
``A Legacy Recalled and a Tradition Continued,''
\emph{Journal of Automated Reasoning}, vol. 27, no. 2, pp. 97--122, 2001.

\bibitem{Meredith1953}
C. A. Meredith,
``A Single Axiom of Positive Logic,''
\emph{Journal of Computing Systems}, vol. 1, no. 3, pp. 169--170, 1953.

\end{thebibliography}

\end{document}
